\section{Finite signed certificates and a complete vertical wall}
\label{sec:finite}

The approximation theorems become exclusion results only after a strict
margin for the complete field has been verified. We describe the finite
error contract, state the interval-certified margins, and prove the
finite-height consequence from the wall and named published inputs.
The computations use outward Arb/Acb arithmetic through
\texttt{python-flint}; the accompanying source and records specify the
library versions, exact cells and frozen parameters. Their status is
computer-assisted interval certification, subject to the written error
bounds and the arithmetic implementation.

\subsection{Retained-source and whole-cell errors}

Let $P_{m,d,U}$ be the paired integral in
\eqref{eq:shifted-finite} with upper endpoint $U$ in place of infinity.
For fixed $m,d,U$ it satisfies $P_t=-P_{zz}$. For each retained
$n<m$ and derivative order $j$, put
\[
 D_{n,j}=4c_d\lambda_ne^{4U}-9-Y-2TU-\frac j{U+d}.
\]
If $D_{n,j}>0$ and $16c_d\lambda_ne^{4U}>2T$, a concave-tangent
bound gives the following retained-source error:
\begin{equation}\label{eq:finite-source-tail}
 e^{-dx}\sum_{n<m}
 \frac{(2\lambda_n^2+3\lambda_n)(U+d)^j
 e^{TU^2+(9+Y)U-c_d\lambda_ne^{4U}}}{D_{n,j}}.
\end{equation}
To see this, take the logarithm of the positive integrand majorant
$(u+d)^j\exp(Tu^2+(9+Y)u-c_d\lambda_ne^{4u})$.
Its derivative at $U$ is $-D_{n,j}$ and its second derivative
is negative on $[U,\infty)$ under the stated condition.
Its tangent therefore integrates to the expression in
\eqref{eq:finite-source-tail}. Sum both contours with their factor $1/2$.

For propagation about $(t_c,x_c)$, expand first in space to degree $L$
and then each coefficient in time to degree $K$. Suppose a common bound
$M$ covers the derivatives of $P_{m,d,U}$ through the needed orders,
divided by a fixed positive normalization $G_*$. For a requested
derivative $j$, those orders include
$j+L+1$ and $j+L+2(K+1)$. If
$|z-x_c|\le R$ and $|t-t_c|\le\Delta$, the normalized remainder is at most
\begin{equation}\label{eq:cell-taylor}
 M\left(\frac{R^{L+1}}{(L+1)!}
       +e^R\frac{\Delta^{K+1}}{(K+1)!}\right).
\end{equation}
Indeed spatial Taylor's theorem gives the first term. The heat equation
replaces each time derivative by a signed even spatial derivative; the
sum of the absolute spatial Taylor weights is at most $e^R$.
The polynomial coefficients are midpoint derivatives enclosed by rigorous
quadrature. Their radii are propagated with the same absolute weights.
Outward evaluation encloses the polynomial on the whole cell.

Finally add the physical omitted-mode and retained-source errors at
the requested orders. Propagation applies to the finite integral itself;
it does not replace its high derivatives by unproved relative errors
for the physical field. The full budget therefore contains mode omission,
source cutoff, quadrature, subsequent rounding and whole-cell variation.

\subsection{Two signed rectangles}

Recall $E=H^2+H_x^2$ and $\Lap_0=H_x^2-HH_{xx}$.
\begin{proposition}[Interval-certified signed margins]\label{prop:finite-signed}
The strict inequality
\begin{equation}\label{eq:finite-signed}
 \Lap_0(t,x)>\frac1{100}E(t,x)
\end{equation}
holds throughout each of the closed rectangles
\[
 D_0=[1/100,1/5]\times[25,45],\qquad
 D_1=[1/100,1/5]\times[999,1001].
\]
\end{proposition}
\begin{proof}[Certificate and error contract]
On $D_0$, the complete-source evaluator certifies the inequality on
314 contiguous full-time spatial strips. The replayed lower excess
$\Lap_0-E/100$ exceeds $4\,10^{-16}$ throughout every strip.
This is a quantitative signed certificate on a previously covered
domain, not a new global sign law.

For $D_1$ freeze
\[
 t_c=21/200,\quad x_c=1000,\quad d=\pi/8-4\log(1000)/1000,
 \quad m=12,\quad U=2,
\]
\[
 G_*=\mathcal G_0(1000)=1000^{7/4}e^{-125\pi},\quad L=24,\quad K=12.
\]
The three physical derivatives require midpoint orders $0$ through $50$
and a common remainder bound through order $52$. One valid normalized
majorant for all these finite-integral derivatives on $D_1$ is
\begin{equation}\label{eq:finite-majorant}
 M=1000^{9/4}e^d\int_0^2
 \sum_{n=1}^{11}(2\lambda_n^2+3\lambda_n)
 e^{u^2/5+9u-c_d\lambda_ne^{4u}}[1+(u+d)^{52}]\dd u.
\end{equation}
The factor $e^d$ pays the left edge $x=999$, and
$(u+d)^j\le1+(u+d)^{52}$ pays every intermediate order.
Outward integration gives $M<9.632889\,10^{10}$.
Formula~\eqref{eq:cell-taylor} pays the entire Taylor remainder.
The mode error uses Theorem~\ref{thm:nonstationary} with the spectral
factor at $x=999$ and the phase gap at $x=1001$; interval checks give
$B>190.59$ and $P>797.37$. Every condition in
\eqref{eq:finite-source-tail} is also checked with intervals.

The closed time slabs have endpoints
\[
 1/100,\quad23/400,\quad21/200,\quad61/400,\quad1/5.
\]
Ten spatial cells in each slab give 40 accepted cells in all.
Exact rational endpoint comparisons verify their full closed coverage.
On each cell, outward evaluation of
$H_x^2-HH_{xx}-(H^2+H_x^2)/100$ has lower bound exceeding
$8.26087\,10^{-5}G_*^2$. The largest added component error is
below $4.621\,10^{-12}G_*$; quadrature radii and polynomial-image
rounding are retained in the enclosures. Square ranges include zero
when their input interval does. These strict comparisons prove
\eqref{eq:finite-signed}.

The numerical supplement retains every cell, the failed coarse choices,
exact dyadic midpoint data, analytic parameters and source identities.
Fresh 256-bit shifted integration and a separately evaluated 768-bit
real-source representation reproduce the high-height conclusions.
These are internal computational checks, not independent specialist review.
\end{proof}

\subsection{A wall covering the complete initial-to-final time interval}

\begin{proposition}[Interval-certified vertical wall]\label{prop:wall}
For $0\le t\le1/5$ and $|y|\le1$,
\begin{equation}\label{eq:wall}
 |H(t,1000+iy)|>0.00425\,G_*.
\end{equation}
\end{proposition}
\begin{proof}[Certificate and error contract]
Reuse the frozen midpoint jets of $D_1$. For the complex segment,
replace the majorant in \eqref{eq:finite-majorant} by
\[
 M_{\rm wall}=1000^{9/4}\int_0^2
 \sum_{n=1}^{11}(2\lambda_n^2+3\lambda_n)
 e^{u^2/5+10u-c_d\lambda_ne^{4u}}[1+(u+d)^{52}]\dd u.
\]
The additional $e^u$ pays $|\Im z|\le1$; the factor $e^d$
is unnecessary because $\Re z=1000$. The largest time distance
from $t_c$ is $21/200$, including $t=0$. Use
\eqref{eq:cell-taylor}, the mode bound with $Y=1$, and
\eqref{eq:finite-source-tail} with $Y=1$.

Seven closed cells cover $[0,1/5]\times[0,1]$ in $(t,y)$,
with time endpoints $0,1/20,1/10,3/20,1/5$.
On each cell, a selected real or imaginary coordinate of $H/G_*$
has absolute lower bound greater than $0.00425150082920$.
Every added coordinate error is below $1.533\,10^{-11}G_*$.
The modulus is at least the modulus of that coordinate. Exact coverage
checks include all seams and endpoints; conjugation supplies $-1\le y\le0$.
Fresh shifted and real-source replays recompute every margin.
\end{proof}

\subsection{Finite propagation without an initial simplicity assumption}

\begin{lemma}[Finite-height propagation]\label{lem:finite-propagation}
Let $F(t,z)$ be analytic in a neighborhood of the relevant time and
spectral domains, satisfy $F_t=-F_{zz}$, and have real conjugation
symmetry for real $t$. Fix $X,T>0$. Suppose, for $0\le t\le T$:
all zeros lie in $|\Im z|\le1$; $F(t,iy)\ne0$ for $|y|\le2$;
and $F(t,X+iy)\ne0$ for $|y|\le1$. Suppose also that all zeros at
$t=0$ with $0\le\Re z\le X$ are real.
Then all zeros in that vertical region remain real for $0\le t\le T$,
and its real zeros are simple for $0<t\le T$.
\end{lemma}
\begin{proof}
The rectangle $\Omega=\{0<\Re z<X,\ |\Im z|<2\}$ has no
boundary zeros: the strip hypothesis pays the horizontal edges and
the upper and lower parts of the right edge; the remaining boundary
conditions pay the rest. The argument principle keeps its finite
zero count, with multiplicity, constant throughout the time interval.
Simple real roots persist real by the implicit function theorem
and conjugation uniqueness.

At a real root of multiplicity $r\ge2$ at $(t_*,x_*)$, the heat
equation and analytic Taylor expansion give the leading weighted term
\[
 a_r(2h)^{r/2}\operatorname{He}_r
       \left(\frac{z-x_*}{\sqrt{2h}}\right),\qquad
 a_r=\frac{\partial_z^rF(t_*,x_*)}{r!},\quad h=t-t_*,
\]
where $\operatorname{He}_r$ is the monic probabilists' Hermite polynomial.
Its roots $\lambda_1<\cdots<\lambda_r$ are real and simple.
The local roots consequently have the expansions
\[
 z_k(t_*+h)=x_*+\sqrt{2h}\lambda_k+O(h).
\]
For $h>0$, small disjoint disks about the real leading roots
each contain exactly one root by Rouch\'e's theorem, and conjugation
makes that root real. For $h<0$, at least one nonzero Hermite
root gives a nonreal predecessor. This is also the local splitting
statement in \cite[Proposition 3.1(ii)]{Polymath}.

The finitely many initial root clusters therefore become real
simple roots for a short positive interval, even if initial roots
are multiple. Compactness, the fixed boundary-free count and
Rouch\'e's theorem make reality closed in time. At a first exit,
a repeated real root would have nonreal predecessors, while a
simple real root would persist real. Neither allows an exit.
A repeated root at any positive time, including $T$, likewise
contradicts reality immediately beforehand. The boundary has no roots.
\end{proof}

\begin{theorem}[Certified finite-height noncollision]\label{thm:finite-height}
In the normalization \eqref{eq:xi},
\begin{align}
 E(t,x)&>0&& (0<t\le1/5,\ |x|\le1000),\label{eq:finite1000}\\
 E(t,x)&>0&& (1/100\le t\le1/5,\ |x|\le1001).\label{eq:finite1001}
\end{align}
\end{theorem}
\begin{proof}
The initial zeta zero strip and de Bruijn's strip-preservation theorem
\cite[Theorem 3.2]{Polymath} keep all zeros of $H(t,\cdot)$
in $|\Im z|\le1$. Positivity of $\Phi$ gives $H(t,iy)>0$
from its hyperbolic-cosine integral. Proposition~\ref{prop:wall}
gives the remaining wall condition.
Platt--Trudgian \cite[Theorem 1]{PlattTrudgian} verify initial
zero reality through zeta height $3\,000\,175\,332\,800$.
Only height $500$ is used here, since \eqref{eq:xi} sends
$\Re z=1000$ to zeta height $500$.
Apply Lemma~\ref{lem:finite-propagation} with $X=1000$, $T=1/5$,
then use evenness. This proves \eqref{eq:finite1000}.
The strict signed inequality on $D_1$ excludes a common zero
of $H,H_x$ there. Its overlap $[999,1000]$ and evenness prove
\eqref{eq:finite1001}. No initial simplicity or zero-velocity
estimate is assumed.
\end{proof}

\subsection{Measured costs and adversarial checks}

The original matched high-height experiment charged all precomputation,
failed attempts and repairs. It took approximately $7.590$ seconds
for the shifted 160-bit calculation and $11.236$ seconds for the
768-bit real-source calculation, a ratio of $1.480$.
Both certified the same 40 signed cells with the same useful
physical derivative majorant. The declared factor-two performance
gate failed. These are recorded bounded timings, not a complexity
theorem or an asserted general advantage of contour shifting.
Separate certified point evaluations at $t=1/10$, $x=1000,10000$
used 15 and 49 retained modes, respectively, and included all
integration errors; point accuracy alone supplies no interval coverage.

The source-specific test interface rejects all five applicable
collision-event boxes in the control catalogue and separately rejects
the Euler-factor negative-Laguerre witness. Finite-prime and
weighted-square sources lack complete normalized analytic reflection;
reflected, primitive and counter-kernel sources lack the unchanged
theta square expansion; Euler-factor controls change the coefficients
and level. They are not assigned theta's truncation constants.
Malformed-cache, missing-cell, changed-source and strict-equality
tests also reject. An exact polynomial heat event tests the same
cell interface with its own zero-truncation contract.
These checks test failure detection and arithmetic implementation;
they do not extend any finite result to unbounded height.
